\documentclass[12pt]{article}
\usepackage[T1]{fontenc}
\usepackage[utf8]{inputenc}
\usepackage{lmodern}
\usepackage{textcomp}
\usepackage{enumitem}
\usepackage{geometry}
\geometry{margin=1in}
\begin{document}
\begin{center}
Answer Key - History - Graduate Level Assessment 6 \
\small (Answers are ordered exactly as in the question paper.)
\end{center}

\section*{Multiple Choice}
\begin{enumerate}[label=\arabic*.]
\item \textbf{Answer:} A
\\ \textit{Options:} A) it expended less energy than quadrupedalism for going long distances; B) it freed the hands to carry things; C) it allowed hominids to see greater distances; D) it provided a more intimidating posture to potential predators; E) it allowed for better heat regulation
\\ \textit{Note:} The statement is correct because while energy efficiency in bipedalism is a widely accepted hypothesis, the assertion that it expended less energy than quadrupedalism for long distances is misleading, as evidence suggests that quadrupedal locomotion is generally more energy-efficient for long-distance travel.
\item \textbf{Answer:} E
\\ \textit{Options:} A) squash; B) knotweed; C) beans; D) sunflower; E) maize
\\ \textit{Note:} Maize was not a significant part of Hopewell subsistence because the Hopewell culture, which thrived from approximately 200 BCE to 500 CE, primarily relied on hunting, gathering, and the cultivation of native plants such as squash and beans, rather than the agricultural practice of maize farming that became prominent later in the region.
\item \textbf{Answer:} B
\\ \textit{Options:} A) Simple foragers employ domestication of animals.; B) Complex foragers focus on a few highly productive resources.; C) Complex foragers employ irrigation technology.; D) Simple foragers use advanced hunting techniques.; E) Complex foragers rely primarily on marine resources.
\\ \textit{Note:} The correct answer highlights that complex foragers employ a more specialized and strategic approach to resource acquisition by concentrating on a limited number of highly productive resources, unlike simple foragers who typically utilize a broader range of less productive ones.
\item \textbf{Answer:} A
\\ \textit{Options:} A) within the past 12,000 years; B) within the past 1,000 years; C) within the past 24,000 years; D) within the past 2,000 years; E) within the past 6,000 years
\\ \textit{Note:} The correct answer is within the past 12,000 years because this timeframe corresponds to the beginning of the Neolithic Revolution, when humans transitioned from nomadic hunter-gatherer lifestyles to settled agricultural practices, actively cultivating and managing food sources.
\item \textbf{Answer:} E
\\ \textit{Options:} A) the Temple of the Feathered Serpent; B) an enormous walled-in precinct of elite residences; C) the Palace of the Painted Walls; D) a complex irrigation system; E) the Pyramid of the Sun
\\ \textit{Note:} The Pyramid of the Sun served as a central architectural and ceremonial structure in Moche society, symbolizing their religious beliefs and social organization while also demonstrating their advanced engineering skills.
\end{enumerate}

\section*{True / False}
\begin{enumerate}[label=\arabic*.]
\setcounter{enumi}{5}
\item \textbf{Answer:} True
\\ \textit{Options:} A) True; B) False
\\ \textit{Note:} According to the passage: German Romanticism
\item \textbf{Answer:} True
\\ \textit{Options:} A) True; B) False
\\ \textit{Note:} According to the passage: chalkos
\item \textbf{Answer:} True
\\ \textit{Options:} A) True; B) False
\\ \textit{Note:} According to the passage: hundreds
\item \textbf{Answer:} True
\\ \textit{Options:} A) True; B) False
\\ \textit{Note:} According to the passage: at the top of the shaft or beside the bottom of the shaft
\item \textbf{Answer:} True
\\ \textit{Options:} A) True; B) False
\\ \textit{Note:} According to the passage: Swainson's thrush
\end{enumerate}

\section*{Long Form Response}
\begin{enumerate}[label=\arabic*.]
\setcounter{enumi}{10}
\item \textbf{Answer:} polyphonic
\\ \textit{Note:} Expected answer: polyphonic
\item \textbf{Answer:} West Africa
\\ \textit{Note:} Expected answer: West Africa
\end{enumerate}

\vfill
\noindent\textit{End of Answer Key}
\end{document}